\ifdefined\gkver\else\def\gkver{student}\fi
\documentclass[\gkver]{gaokaozhenti}
\begin{document}

\chapter{2001年广东综合}
%% paperType: 综合能力测试物理部分

\begin{enumerate}
\item
%% number: 19
%% paperName: 2001年广东综合第 19 题 3 分
%% typeId: 1
%% type: 单选题
%% chapter: 原子和原子核
%% point: 天然放射性
%% method: 无
%% score: 3
%% degree: 850
%% duplicateId: 0
%% body:
原子序数大于 92 的所有元素，都能自发地放出射线，这些射线共有三种：$\alpha$ 射线、$\beta$ 射线和 $\gamma$ 射线。下列说法中正确的是\xzanswer{D}

\fourchoices[answer=D]
{原子核每放出一个 $\alpha$ 粒子，原子序数减少 4}
{原子核每放出一个 $\alpha$ 粒子，原子序数增加 4}
{原子核每放出一个 $\beta$ 粒子，原子序数减少 1}
{原子核每放出一个 $\beta$ 粒子，原子序数增加 1}

%% answer: D
%% memo:
\memoanswer{
本题原卷及材料中未提供详解，待补充。
}

\item
%% number: 20
%% paperName: 2001年广东综合第 20 题 3 分
%% typeId: 1
%% type: 单选题
%% chapter: 原子和原子核
%% point: 核裂变
%% method: 无
%% score: 3
%% degree: 850
%% duplicateId: 0
%% body:
原子核反应有广泛的应用，如用于核电站等。在下列核反应中，属于核裂变反应的是\xzanswer{D}

\fourchoices[answer=D]
{\ce{^{10}_{5}B + ^{1}_{0}n -> ^{7}_{3}Li + ^{4}_{2}He}}
{\ce{^{238}_{92}U -> ^{234}_{90}Th + ^{4}_{2}He}}
{\ce{^{14}_{7}N + ^{4}_{2}He -> ^{17}_{8}O + ^{1}_{1}H}}
{\ce{^{235}_{92}U + ^{1}_{0}n -> ^{141}_{56}Ba + ^{92}_{36}Kr + 3^{1}_{0}n}}

%% answer: D
%% memo:
\memoanswer{
本题原卷及材料中未提供详解，待补充。
}

\item
%% number: 28
%% paperName: 2001年广东综合第 28 题 3 分
%% typeId: 1
%% type: 单选题
%% chapter: 电磁感应
%% point: 楞次定律
%% method: 无
%% score: 3
%% degree: 650
%% duplicateId: 0
%% body:
有一种高速磁悬浮列车的设计方案是在每节车厢底部安装强磁铁（磁场方向向下），并在两条铁轨之间沿途平放一系列线圈。下列说法中不正确的是\xzanswer{D}

\fourchoices[answer=D]
{当列车运动时，通过线圈的磁通量会发生变化}
{列车速度越快，通过线圈的磁通量变化越快}
{列车运动时，线圈中会产生感应电流}
{线圈中的感应电流的大小与列车速度无关}

%% answer: D
%% memo:
\memoanswer{
本题原卷及材料中未提供详解，待补充。
}

\item
%% number: 29
%% paperName: 2001年广东综合第 29 题 3 分
%% typeId: 1
%% type: 单选题
%% chapter: 力物体的平衡
%% point: 摩擦力
%% method: 无
%% score: 3
%% degree: 850
%% duplicateId: 0
%% body:
磁悬浮列车在行进时会“浮”在轨道上方，从而可高速行驶。可高速行驶的原因是：列车浮起后\xzanswer{C}

\fourchoices[answer=C]
{减少了列车的惯性}
{减少了地球对列车的引力}
{减少了列车与铁轨间的摩擦力}
{减少了列车所受的空气阻力}

%% answer: C
%% memo:
\memoanswer{
本题原卷及材料中未提供详解，待补充。
}

\item
%% number: 30
%% paperName: 2001年广东综合第 30 题 3 分
%% typeId: 1
%% type: 单选题
%% chapter: 功和能
%% point: 功率
%% method: 无
%% score: 3
%% degree: 750
%% duplicateId: 0
%% body:
假设列车从静止开始匀加速运动，经过 $500\Um$ 的路程后，速度达到 $360\Ukmh$。整个列车的质量为 $1.00\times10^{5}\Ukg$，如果不计阻力，在匀加速阶段、牵引力的最大功率是\xzanswer{B}

\fourchoices[answer=B]
{$4.67\times10^{6}\UkW$}
{$1.0\times10^{5}\UkW$}
{$1.0\times10^{8}\UkW$}
{$4.67\times10^{9}\UkW$}

%% answer: B
%% memo:
\memoanswer{
本题原卷及材料中未提供详解，待补充。
}

\item
%% number: 33
%% paperName: 2001年广东综合第 33 题 4 分
%% typeId: 3
%% type: 填空题
%% chapter: 功和能
%% point: 功
%% method: 无
%% score: 4
%% degree: 750
%% duplicateId: 0
%% body:
用塑料部件代替金属部件可以减小汽车质量，从而可以节能。假设汽车使用塑料部件后质量为原来的 $4/5$，而其它性能不变。若车行驶时所受的阻力与车的质量成正比，则与原来相比，使用塑料部件的汽车能节省\tkanswer[1.5cm]{20}\%的汽油。

%% answer: 20
%% memo:
\memoanswer{
本题原卷及材料中未提供详解，待补充。
}

\item
%% number: 38
%% paperName: 2001年广东综合第 38 题 4 分
%% typeId: 6
%% type: 计算题
%% chapter: 功和能
%% point: 功率
%% method: 无
%% score: 4
%% degree: 650
%% duplicateId: 0
%% body:
某市计划每日供水 $180\times10^{4}\Ut$，在市郊修建了一水库。为了将水送入水库，需要将水渠的水提高 $30\Um$。设每根输水管水泵功率为 $100\UkW$，且水泵昼夜不停地工作。如不计机械能的损耗，至少需要安装多少根输水管？每根输水管中每秒流过的水量为多少吨？（取 $g=10\Umsq$。）

\begin{enumerate}
	\item
	至少需要安装多少根输水管？
	\item
	每根输水管中每秒流过的水量为多少吨？
\end{enumerate}

%% answer: 63 根输水管。0.33 吨
\jdanswer{
\begin{enumerate}
	\item
	63 根输水管。
	\item
	$0.33\Ut$。
\end{enumerate}
}

%% memo:
\memoanswer{
将 $180\times10^{4}\Ut$ 水提高 $30\Um$ 需做的功为

$W = mgh = 180 \times 10^{4} \times 10^{3} \times 10 \times 30\UJ$

（1 分）

每台水泵每昼夜所做的功为

$W_{0} = Pt = 100\times10^{3}\times24\times3600\UJ$

（1 分）

两者相除得到

$\frac{W}{W_{0}} = 62.5$

（1 分）

由于每台水泵配一根输水管，故至少需要 63 根输水管。

每秒流过一根水管的水量为

$M = 180 \times 10^{4} / (63 \times 24 \times 3600) = 0.33\Ut$
}

\end{enumerate}

\end{document}
